\section{引言}
\label{sec:introduction}
空地协作将执行观测的机器人与最终执行操作的机器人分离。
无人机能够从多个视点检查作业区域；地面移动操作机器人则能够取回检测到的物体。
然而，成功协作不仅要求从空中检测目标，还要求地面机器人占据适当站位、在近距离看清目标，
并完成抓取与提升。因此，感知决策不仅是寻找揭示最多未知空间的视点，
更是寻找能够补足后续操作所需证据的视点。

下一最佳视点（next-best-view，NBV）规划为选择观测位置提供了自然工具。
探索和重建方法通常根据尚未建图的空间或观测不完整的表面评价测量价值
\cite{bircher2016nbv,border2018see}。这些目标适合其原有建图任务，
但本身并未表达另一台机器人随后执行操作的能力。
对于取回任务，远离有效机械臂站位的大面积未知区域可能具有较高的通用观测增益，
却不能提供可执行操作位姿所需的证据。相反，一个与高价值底盘占地轮廓相交的小型未观测区域，
可能阻止整个任务继续。在有限感知预算下，观测的空间分配由此成为关键。

面向任务的视点选择并非新概念。已有研究针对抓取检测选择视点，
或联合决定继续观测与开始抓取的时机\cite{gualtieri2017viewpoint,breyer2022target}。
本文研究不同的耦合关系：空中传感器并不直接执行抓取，所需观测的区域也不限于物体表面。
它还取决于地面机械臂可以站在哪里、底盘占据多大范围，以及从站位到完成取回之间的操作约束。
本文的研究问题是：在相同观测预算下，这些下游操作需求能否比信息中心型 NBV 目标
更有效地引导空中观测？

核心思路是将未来操作需求反向传递给感知目标。经验证的逆可达性候选确定潜在地面站位；
精确站位的占地轮廓将操作相关性映射到环境空间，强调有用站位附近尚未解决的观测缺失。
有限扫描模型预测下一视点的采集机会，而共用执行筛选器判断实测证据是否支持任务交接。
这一表示将视点选择连接到物理取回，而非以地图覆盖率或单个可达抓取作为终点。

本文的贡献包括：
\begin{enumerate}
\item 提出面向异构空地取回任务的操作需求引导主动感知框架，使地面机器人的未来操作需求
指导无人机的感知决策。
\item 提出考虑底盘占地轮廓的任务相关性表示，将经验证的机械臂可达性、精确底盘位姿和
操作几何与环境观测缺失度结合，并区分未评估的操作能力、未观测地面和已观测碰撞证据。
\item 在 P450、BUNKER、AUBO i5 和 AG95 物理仿真系统中完成 96 个独立场景、
192 次主配对任务的闭环验证，以物理取回成功为主要终点。
Generic/Ours 共用传感器模型、候选、成本、预算与执行链，隔离观测增益加权这一差异。
\end{enumerate}
本文的贡献不包括新的逆运动学求解器，也不将共用平台修复计为方法优势。
成功率结论限于定义的仿真任务总体；效率、失败机制与迁移边界分别讨论。
